The results indicate that oil prices are strongly influenced by global economic and political events.
\subsection{1973 Oil Crisis}
On October 19,1973 the Organization of Arab Petroleum Exporting Countries also called (OAPEC) issued an oil embargo on the United States. This was issued due to president Nixon's request to congress to aid Israel during the conflict of Yom Kippur War. The embargo ceased U.S. oil imports from participating in OAPEC. Also, production cute began to take place which ultimately altered oil prices around the globe. Based on the data we can see oil in January of 1974 rose from 2.90 a barrel to 11.65 per barrel in March \cite{oil1973}.
\subsection{1978 Oil Shock}
The 1978 Oil Shock is closely related to the 1973 Oil Crisis. However the increase in oil prices was much larger compared to the 1973 crisis. Pertoleum oil grew from 15.85 to 39.50 per barrel. The origin of this crisis is based of the instability in Iran. In 1978 we see a widespread strike by oil workers that greatly reduced the production of oil. Meanwhile, Irans government was undergoing a political change, where the departure of Shah Mohammad Reza Pahlavi led to the cessation of Iranian oil exports. When looking at the effects of this change we see that global oil supply reduced by 4-5 percent. However due to the previous crisis the impact was far more substantial than before. This caused countries all around the globe to go into panic especially the United States. We saw people panic buying, distrubution issues, and a shortage at the pump. This ultimately caused stagflation where costs kept rising compared to the slow economic growth \cite{oil1979}.
\subsection{Crash of Oil Prices in 1986}
This is different compared to the other two events. We see oil prices drop by more than 60 percent within a year. This occurred due to the OPEC embargo finally ending causing countries to be able to purchase oil from other countries. As a result, this led to the United states to have an excess amount of oil. Due to this we see Oil fields to begin shutting down and many workers losing their jobs. As a result, the United states shifted from making their own oil to buying their oil for cheaper from overseas. This event highlights that even if oil prices decrease it doesn't connect to economic growth increasing \cite{oil1986}.
\subsection{OPEC Production Cut}
This event highlights the greed behind making more money off the population. In 1999 OPEC decides to cut oil production 4 million barrels of oil. The  reason for this is because OPEC wants to raise the price of oil and many countries agreed to further cut oil production in order to increase prices. The goal of this was to increase prices from 10 to 16-17 per barrell \cite{opec1999}.
\subsection{2008 Financial Crisis}
The 2008 Financial Crisis was the effect of the 2006 housing market crash. The effects of the housing market crash began to spread into the Oil industry in June of 2008. Oil prices were at 139 and fell to 39 in feburary of 2009. The decrease in price was due to the financial crisis. This occured due to the lack of demand and the lack of money that the population had. As a result, the population was unable to afford the prices at the pump causing the prices to fall drastically \cite{crisis2008}.
\subsection{2016 Oil Crisis}
After the rise of oil prices after the 2008 Financial Crisis disappeared. The economy began to see a rapid shift of oil reducing drastically from 100 in July of 2014 to 30 by Feburary 2016. Unlike other crisis this had long lasting impacts on the worlds economy. What makes this crisis an outlier is that it did not occur during a major political time or due to a political policy being created \cite{imf2017}.
\subsection{COVID-19 Pandemic}
This event resulted in the largest shift in Oil prices since the late 1900's. When COVID-19 was first detected in China we saw the demand for oil in China decrease by 3 million barrels per day. What makes this unique compared to other events, is that main producers of oil thought of this as a situation to produce excess amounts of oil to gain power and strength. However, this did not occur, due to the excess amount of oil and the lack of demand we see countries like the united states have 535.2 million barrels of oil in stockpile. The oil prices continued to decrease more and more due to everyone being isolated and not needing to travel. As a result, we see that oil fell over 60 percent in 4 months. The oil prices began to increase when OPEC decided to meet and cut oil production in order to reduce the surplus. Oil prices began to slowly recover as travel restrictions decreased and COVID-19 slowing down \cite{covidoil}.
